% ============================================================================
%  IV.5 Équations de Gibbs-Appell
% ============================================================================
\subsection{Équations de Gibbs-Appell}

La formulation de Gibbs-Appell écrit la dynamique à partir des
accélérations des points matériels, en utilisant directement les
vitesses indépendantes autorisées par les contraintes. On se place
dans un référentiel galiléen, avec des masses constantes et des
liaisons idéales, et des contraintes de la forme
\[
    \mat{A}(q)\,\dot q = 0
\]

% ----------------------------------------------------------------------------
\subsubsection{Paramétrer les vitesses admissibles}

Soit $q\in\R^n$ un système de coordonnées locales, après prise en
compte des contraintes holonomes. On suppose
\[
    \operatorname{rang}\mat{A}(q) = k, \qquad d = n - k
\]

\begin{definition}[Matrice des vitesses admissibles]
    On choisit une matrice $\mat{N}(q)\in\mathcal{M}_{n,d}(\R)$ dont
    les colonnes $\vect{N}_1,\dots,\vect{N}_d$ forment une base de
    $\ker\mat{A}(q)$ telle que : 
    \[
        \mat{A}(q)\,\mat{N}(q) = 0, \qquad
        \operatorname{rang}\mat{N}(q) = d
    \]
\end{definition}

\begin{proposition}[Pseudo-vitesses]
    Toute vitesse admissible s'écrit de manière unique
    \[
        \dot q = \mat{N}(q)\,u = \sum_{j=1}^{d} u_j\,\vect{N}_j(q),
        \qquad u\in\R^d
    \]
    où $u$ est le vecteur des \textbf{pseudo-vitesses}. De même, les
    déplacements virtuels admissibles s'écrivent
    \[
        \delta q = \mat{N}(q)\,\xi, \qquad \xi\in\R^d
    \]
\end{proposition}
Les pseudo-vitesses sont les vitesses libres du système sous
contraintes. En fait, au lieu de $n$ vitesses généralisées liées par $k$
équations, on ne manipule que $d = n - k$ vitesses indépendantes, et
la contrainte $\mat{A}(q)\,\dot q = 0$ se vérifie automatiquement.
La dynamique se réduit alors à $d$ équations, sans multiplicateurs de
Lagrange.

% ----------------------------------------------------------------------------
\subsubsection{Accélérations et fonction de Gibbs-Appell}

On considère des points matériels de masses $m_a$ et de positions
$\vect{r}_a = \vect{r}_a(q)$. Leurs vitesses sont
\[
    \vect{v}_a = \frac{\partial\vect{r}_a}{\partial q}\,\dot q
    = \frac{\partial\vect{r}_a}{\partial q}\,\mat{N}(q)\,u
\]
On pose
\[
    \mat{B}_a(q) = \frac{\partial\vect{r}_a}{\partial q}\,\mat{N}(q)
    \in\mathcal{M}_{3,d}(\R),
    \qquad\text{d'où}\qquad
    \vect{v}_a = \mat{B}_a(q)\,u
\]
On dérive le long de la trajectoire (on se rappelle que $\vect{v}_a$ ne dépend pas explicitement de $t$) :
\[
    \vect{a}_a = \frac{\mathrm{d}\vect{v}_a}{\mathrm{d}t}
    = \mat{B}_a(q)\,\dot u + \dot{\mat{B}}_a\,u,
    \qquad
    \dot{\mat{B}}_a = \frac{\partial\mat{B}_a}{\partial q}\,\dot q
\]
où $\dot{\mat{B}}_a\,u$ ne dépend que de $q$ et $u$.

\begin{proposition}[Accélérations]
    L'accélération $\vect{a}_a$ dépend de $\dot u$ de façon affine, et
    \[
        \frac{\partial\vect{a}_a}{\partial\dot u_j}
        = \vect{B}_{a,j}
        = \frac{\partial\vect{v}_a}{\partial u_j},
        \qquad j = 1,\dots,d
    \]
    où $\vect{B}_{a,j}$ est la $j$-ième colonne de $\mat{B}_a$.
\end{proposition}

\begin{definition}[Fonction de Gibbs-Appell]
    La \textbf{fonction de Gibbs-Appell}, ou énergie des
    accélérations, est
    \[
        \mathcal{G}(q, u, \dot u)
        = \frac{1}{2}\sum_a m_a\,\|\vect{a}_a\|^2
    \]
\end{definition}

% ----------------------------------------------------------------------------
\subsubsection{Forces associées aux pseudo-vitesses}

Soient $\vect{F}_a$ les forces appliquées. Pour un déplacement virtuel
admissible $\delta q = \mat{N}(q)\,\xi$,
\[
    \delta\vect{r}_a = \pd{\vect{r}_a}{q}\,\delta q
    = \mat{B}_a(q)\,\xi
\]
et leur travail virtuel vaut
\[
    \delta W = \sum_a\vect{F}_a\cdot\bigl(\mat{B}_a(q)\,\xi\bigr)
    = \sum_{j=1}^{d}\mathcal{Q}_j\,\xi_j
\]

\begin{definition}[Forces associées aux pseudo-vitesses]
    La force associée à la pseudo-vitesse $u_j$ est le coefficient
    $\mathcal{Q}_j$ de $\xi_j$ dans le travail virtuel :
    \[
        \mathcal{Q}_j = \sum_a\vect{F}_a\cdot\vect{B}_{a,j}
        = \sum_a\vect{F}_a\cdot\pd{\vect{v}_a}{u_j}
    \]
\end{definition}

\begin{proposition}[Lien avec les forces généralisées]
    Si $\vect{Q}$ désigne les forces généralisées associées aux
    coordonnées $q$, alors
    \[
        \mathcal{Q} = \mat{N}(q)^T\vect{Q}
        = \mat{N}(q)^T\left(-\pd{V}{q} + \vect{Q}^{\mathrm{nc}}\right)
    \]
\end{proposition}

\begin{demonstration}
    Comme $\mat{B}_a (q)= \mat{J}_a(q)\mat{N}(q)$,
    $\delta W = \sum_a\vect{F}_a^T\mat{J}_a\mat{N}\,\xi
    = \vect{Q}^T\mat{N}\,\xi$, et on identifie les coefficients de
    $\xi$.
\end{demonstration}

% ----------------------------------------------------------------------------
\subsubsection{Théorème de Gibbs-Appell}

\begin{theoreme}[Gibbs-Appell]
    Sous les hypothèses précédentes, les équations du mouvement sont
    \[
        \pd{\mathcal{G}}{\dot u_j} = \mathcal{Q}_j,
        \qquad j = 1,\dots,d
    \]
    complétées par la relation cinématique
    \[
        \dot q = \mat{N}(q)\,u
    \]
\end{theoreme}

\begin{demonstration}
    Les liaisons sont idéales : d'après le principe de d'Alembert, pour
    tout déplacement virtuel admissible $\delta\vect{r}_a =
    \mat{B}_a\,\xi$,
    \[
        \sum_a(\vect{F}_a - m_a\vect{a}_a)\cdot\mat{B}_a\,\xi = 0
        \quad\longleftrightarrow\quad
        \sum_{j=1}^{d}\Bigl(\mathcal{Q}_j
        - \sum_a m_a\,\vect{a}_a\cdot\vect{B}_{a,j}\Bigr)\,\xi_j = 0
    \]
    Comme $\xi\in\R^d$ est arbitraire, chaque coefficient est nul :
    \[
        \mathcal{Q}_j = \sum_a m_a\,\vect{a}_a\cdot\vect{B}_{a,j}
    \]
    Or, d'après la proposition sur les accélérations,
    \[
        \pd{\mathcal{G}}{\dot u_j}
        = \sum_a m_a\,\vect{a}_a\cdot\pd{\vect{a}_a}{\dot u_j}
        = \sum_a m_a\,\vect{a}_a\cdot\vect{B}_{a,j}
    \]
    d'où le résultat.
\end{demonstration}

% ----------------------------------------------------------------------------
\subsubsection{Fonction de Gibbs-Appell d'un solide}

Pour un solide $\mathcal{B}$ de masse $m$ et de centre de masse $C$,
la somme sur les points devient une intégrale :
\[
    \mathcal{G} = \frac{1}{2}\int_{\mathcal{B}}\|\vect{a}_P\|^2\,\dd m
\]
avec, d'après le champ des accélérations et en notant
$\vect{r} = \overrightarrow{CP}$,
\[
    \vect{a}_P = \vect{a}_C + \dot{\vect{\Omega}}\wedge\vect{r}
    + \vect{\Omega}\wedge(\vect{\Omega}\wedge\vect{r})
\]

\begin{proposition}[Fonction de Gibbs-Appell d'un solide]
    À des termes indépendants de $\dot{\vect{\Omega}}$ et de
    $\vect{a}_C$ près,
    \[
        \mathcal{G} = \frac{1}{2}\,m\,\|\vect{a}_C\|^2
        + \frac{1}{2}\,\dot{\vect{\Omega}}^T\mat{I}_C\,\dot{\vect{\Omega}}
        + \dot{\vect{\Omega}}\cdot\bigl(\vect{\Omega}\wedge
        \mat{I}_C\vect{\Omega}\bigr)
    \]
    soit, dans le repère du corps,
    \[
        \mathcal{G} = \frac{1}{2}\,m\,\|\vect{a}_C\|^2
        + \frac{1}{2}\,\dot{\vect{\omega}}^T\mat{I}_b\,\dot{\vect{\omega}}
        + \dot{\vect{\omega}}\cdot\bigl(\vect{\omega}\wedge
        \mat{I}_b\vect{\omega}\bigr)
    \]
\end{proposition}

\begin{demonstration}
    On développe le carré de l'accélération :
    \[
        \|\vect{a}_P\|^2
        = \|\vect{a}_C\|^2
        + 2\,\vect{a}_C\cdot\bigl(\dot{\vect{\Omega}}\wedge\vect{r}
        + \vect{\Omega}\wedge(\vect{\Omega}\wedge\vect{r})\bigr)
        + \|\dot{\vect{\Omega}}\wedge\vect{r}\|^2
        + 2\,(\dot{\vect{\Omega}}\wedge\vect{r})\cdot
        \bigl(\vect{\Omega}\wedge(\vect{\Omega}\wedge\vect{r})\bigr)
        + \|\vect{\Omega}\wedge(\vect{\Omega}\wedge\vect{r})\|^2
    \]
    On intègre chaque terme sur $\mathcal{B}$.

    Comme $C$ est le centre de masse,
    $\int_{\mathcal{B}}\vect{r}\,\dd m = 0$, et le terme croisé avec
    $\vect{a}_C$, linéaire en $\vect{r}$, s'annule :
    \[
        \frac12\int_{\mathcal{B}}\|\vect{a}_C\|^2\,\dd m
        + \int_{\mathcal{B}}\vect{a}_C\cdot\bigl(\dot{\vect{\Omega}}
        \wedge\vect{r} + \vect{\Omega}\wedge(\vect{\Omega}\wedge
        \vect{r})\bigr)\,\dd m
        = \frac12\,m\,\|\vect{a}_C\|^2
    \]

    Concernant le terme quadratique en $\dot{\vect{\Omega}}$, on a :
    \[
        \|\dot{\vect{\Omega}}\wedge\vect{r}\|^2
        = \dot{\vect{\Omega}}^T\bigl(\|\vect{r}\|^2\Id
        - \vect{r}\vect{r}^T\bigr)\dot{\vect{\Omega}}
        \quad\longrightarrow\quad
        \frac12\int_{\mathcal{B}}\|\dot{\vect{\Omega}}\wedge\vect{r}\|^2
        \,\dd m
        = \frac12\,\dot{\vect{\Omega}}^T\mat{I}_C\,\dot{\vect{\Omega}}
    \]

    On a
    \[
        \vect{\Omega}\wedge(\vect{\Omega}\wedge\vect{r})
        = (\vect{\Omega}\cdot\vect{r})\,\vect{\Omega}
        - \|\vect{\Omega}\|^2\,\vect{r}
    \]
    Comme $\dot{\vect{\Omega}}\wedge\vect{r}$ est orthogonal à
    $\vect{r}$, seul le premier reste :
    \[
        (\dot{\vect{\Omega}}\wedge\vect{r})\cdot
        \bigl(\vect{\Omega}\wedge(\vect{\Omega}\wedge\vect{r})\bigr)
        = (\vect{\Omega}\cdot\vect{r})\,
        (\dot{\vect{\Omega}}\wedge\vect{r})\cdot\vect{\Omega}
        = (\vect{\Omega}\cdot\vect{r})\,
        \dot{\vect{\Omega}}\cdot(\vect{r}\wedge\vect{\Omega})
    \]
    Or 
    \[
        \mat{I}_C\vect{\Omega}
        = \int_{\mathcal{B}}\bigl(\|\vect{r}\|^2\vect{\Omega}
        - (\vect{r}\cdot\vect{\Omega})\,\vect{r}\bigr)\,\dd m
        \quad\longrightarrow\quad
        \vect{\Omega}\wedge\mat{I}_C\vect{\Omega}
        = \int_{\mathcal{B}}(\vect{\Omega}\cdot\vect{r})\,
        (\vect{r}\wedge\vect{\Omega})\,\dd m
    \]
    d'où
    \[
        \int_{\mathcal{B}}(\dot{\vect{\Omega}}\wedge\vect{r})\cdot
        \bigl(\vect{\Omega}\wedge(\vect{\Omega}\wedge\vect{r})\bigr)
        \,\dd m
        = \dot{\vect{\Omega}}\cdot\bigl(\vect{\Omega}\wedge
        \mat{I}_C\vect{\Omega}\bigr)
    \]

    Finalement, pour le terme
    $\|\vect{\Omega}\wedge(\vect{\Omega}\wedge\vect{r})\|^2$, il ne dépend
    ni de $\dot{\vect{\Omega}}$ ni de $\vect{a}_C$, donc on n'a aucun intérêt de le considérer, on le fait donc disparaître de notre équation.

    \textbf{Forme dans le corps : } Avec $\vect{\Omega} =
    \mat{R}\vect{\omega}$ et $\dot{\mat{R}}\vect{\omega} =
    \mat{R}\,(\vect{\omega}\wedge\vect{\omega}) = 0$,
    \[
        \dot{\vect{\Omega}} = \mat{R}\,\dot{\vect{\omega}},
        \qquad
        \mat{I}_C = \mat{R}\,\mat{I}_b\,\mat{R}^T
    \]
    donc
    \[
        \dot{\vect{\Omega}}^T\mat{I}_C\,\dot{\vect{\Omega}}
        = \dot{\vect{\omega}}^T\mat{I}_b\,\dot{\vect{\omega}},
        \qquad
        \dot{\vect{\Omega}}\cdot(\vect{\Omega}\wedge\mat{I}_C\vect{\Omega})
        = \dot{\vect{\omega}}\cdot(\vect{\omega}\wedge\mat{I}_b\vect{\omega})
    \]
    car $\mat{R}$ conserve le produit scalaire et le produit vectoriel.
\end{demonstration}